\chapter{Template de Capitulo}
\label{chap:template_capitulo}

Este apendice serve como guia de expansao iterativa para novos capitulos.

\section{Objetivos do capitulo}
\begin{itemize}[leftmargin=*]
  \item Objetivo 1
  \item Objetivo 2
\end{itemize}

\section{Implementacao pratica}
\begin{enumerate}[leftmargin=*]
  \item Passo 1
  \item Passo 2
  \item Passo 3
\end{enumerate}

\section{Politica de autocontencao (obrigatoria)}
\begin{itemize}[leftmargin=*]
  \item Todo codigo citado no capitulo deve existir no projeto com caminho valido.
  \item Cada comando deve executar sem dependencias impl\'icitas fora do reposit\'orio.
  \item Sempre que houver \texttt{Plano de testes}, incluir evidencias de execucao no capitulo
        ou referenciar artefatos no Ap\^endice~\ref{apendice:codigo_completo}.
  \item Antes da versao final, executar: \texttt{pwsh -File livro/validar\_autocontencao.ps1}.
\end{itemize}

\section{Plano de testes}
\begin{itemize}[leftmargin=*]
  \item Teste funcional
  \item Teste de performance
  \item Teste de qualidade de resposta
\end{itemize}

\section{Criterios de aceite}
\begin{itemize}[leftmargin=*]
  \item Criterio 1
  \item Criterio 2
\end{itemize}
